\section{Actual recursively constructed lower values along original short histories}

The invariant is the actual recursively constructed family $\operatorname{invFamily}(n,k)$.
No property of an arbitrary lower invariant, desired value equality or
boundary-insertion certificate is assumed.  The initial marked ideal and
entire actual original histories remain fixed.  The proved original short
positive, short zero-contact and crossing diagrams provide the selected
records; here the original short diagrams are consumed by actual values.
The crossing diagram has a different actual lower tail and still requires
recursive transport through that history.

\begin{theorem}[The whole actual zero lower value, in every recursive dimension]
\label{hir:original-interior-full-family-zero-lower-value}
\lean{Hironaka.HironakaOriginalShortLowerValues.original_zero_lowerValue}
\leanok
\source{wlodarczyk05-hironaka,bierstone-milman08-functoriality}
\uses{hir:original-interior-actual-zero-lower-record,hir:original-invariant-empty-actual,hir:step-one-value-masked}
Let $n>0$, $k\geq0$, $M$ originally well formed, $H$ an actual test history
of $\overline M$, $P$ a nonempty actual tail and $Q$ any subsequent actual
history.  Set $j=|M.E|+|H|$.  The actual zero lower value of $M$ formed at
$H\epsilon P$ and computed along $Q$ is
\[
 \operatorname{WithTop.map}(\operatorname{shiftFamily}(k,j))
    \bigl(\operatorname{lowerValue}_k(M,0,HP,Q,y)\bigr).
\]
Both lower values use the unchanged original family and the unchanged
original lowerMarkedM definition, including the actual missing-mask top value.
\end{theorem}
\begin{proof}
The actual whole zero lower record theorem supplies its future empty
insertion at $j$, the whole mask equality and the actual lower-boundary
index bound.  Identify the equal actual opens, keeping their genuine
dependent carrier types.  If the projection of $y$ misses the old mask,
it misses the new equal mask: both values are top, which the map preserves.
Otherwise the real lower value is the actual invariant on the real
restricted $Q$ at the corresponding actual subtype point.  The inserted
lower record and the proved actual all-$k$ invariant insertion theorem
give precisely shiftFamily$(k,j)$ of the old value on that same whole
restricted history.  Coercing into WithTop gives the stated formula.
The needed insertion theorem is proved for this specific recursive family
and supplied internally; it is not a hypothesis about a proposed invariant.
\end{proof}

\begin{theorem}[The positive lower value at its actual index]
\label{hir:original-interior-full-family-positive-lower-value}
\lean{Hironaka.HironakaOriginalShortLowerValues.original_positive_lowerValue}
\leanok
\source{wlodarczyk05-hironaka,bierstone-milman08-functoriality}
\uses{hir:original-interior-actual-positive-lower-record,hir:original-interior-actual-positive-lower-index,hir:original-invariant-empty-actual}
With the same original inputs, at every positive lower index $s+1$ put
\[
 p=Q.\operatorname{proj}(y),
 \qquad c=\operatorname{interiorPositiveLowerIndex}(M,H,P,p).
\]
The actual lower value formed at $H\epsilon P$ is
WithTop.map$(\operatorname{shiftFamily}(k,c))$ of the actual lower value
formed at $HP$, both computed along the same literal $Q$.
\end{theorem}
\begin{proof}
The true member filter computes $c$ in the actual whole lower boundary,
and the actual positive lower producer supplies the whole record insertion
there and equality of the actual masks.  The previous argument now applies
with this actual $c$: the missing-mask branch is top on both sides, while
the genuine masked branch invokes the proved specific all-$k$ insertion
theorem on the actual lower record and the entire actual restriction of
$Q$.  No empty-prefix hypothesis is needed.  The index may lie beyond the
empty original prefix, so no uniform replacement of $c$ by zero is asserted.
\end{proof}

\begin{theorem}[The literal own-level short zero Step One value]
\label{hir:original-interior-full-family-short-zero-step-one-value}
\lean{Hironaka.HironakaOriginalShortLowerValues.actual_short_zero_stepOne_value}
\leanok
\source{wlodarczyk05-hironaka,bierstone-milman08-functoriality}
\uses{hir:original-interior-selected-all-level-short-splits,hir:original-interior-selected-zero-short-lower-diagram,hir:original-interior-full-family-zero-lower-value,hir:original-interior-empty-step-one-dictionary}
Let $S$ be an actual tail and $y$ an actual point whose original current
level is zero and whose actual zero-level run $r$ satisfies $r<|S|$.
For every $k$, the actual Step One value over invFamily$(n,k)$ along
$H\epsilon S$ equals
\[
 \operatorname{pairShift}(\operatorname{shiftFamily}(k,j))
    \bigl(\operatorname{stepOneInvM}_k(M,HS,y)\bigr),
 \qquad j=|M.E|+|H|.
\]
No current nonspecialness or canonical-lower-value certificate is required.
\end{theorem}
\begin{proof}
The actual level and special-point dictionaries identify the corresponding
original branches.  At a special point both actual values are $(0,\top)$,
and the map preserves this value.  Otherwise the actual short split theorem
computes the new current-level run as $r$ and supplies both complete split
records.  Their actual heads are $HP$ and $H\epsilon P$, with the same
literal $Q$, where $P=\operatorname{head}(S,r)$ and $Q=\operatorname{tail}(S,r)$.
The strict run inequality gives $|P|=|S|-r>0$ internally.  Apply the first
theorem to these actual inputs; the lower value is changed by exactly
WithTop.map$(\operatorname{shiftFamily}(k,j))$.  The unchanged outer level
zero gives the asserted whole lexicographic pair formula.
\end{proof}

\begin{theorem}[The literal own-level short positive Step One value]
\label{hir:original-interior-full-family-short-positive-step-one-value}
\lean{Hironaka.HironakaOriginalShortLowerValues.actual_short_positive_stepOne_value}
\leanok
\source{wlodarczyk05-hironaka,bierstone-milman08-functoriality}
\uses{hir:original-interior-selected-positive-short-splits,hir:original-interior-selected-positive-short-lower-diagram,hir:original-interior-full-family-positive-lower-value,hir:original-interior-empty-step-one-dictionary}
If instead the actual current level is $s+1$ and its actual run $r$ is
shorter than $|S|$, put $P=\operatorname{head}(S,r)$,
$Q=\operatorname{tail}(S,r)$, $p=Q.\operatorname{proj}(y)$ and let $c$
be the actual filtered lower insertion index.  The actual new Step One
value is pairShift$(\operatorname{shiftFamily}(k,c))$ of the actual old value.
\end{theorem}
\begin{proof}
Both actual levels and special branches agree.  At special points the
lower entry is top and is preserved.  At nonspecial points the proved
positive short splits supply the identical actual run, selected stage and
tail $Q$.  The computed selected prefix has positive length by the strict
run inequality.  The second theorem gives the actual lower-value map at
its genuine computed index $c$, and retaining the outer actual level
$s+1$ gives the pair formula.  No comparison of different computed indices
or both-zero monomial spectra is inferred from this pointwise equation.
\end{proof}

\begin{theorem}[Order and dominance on the entire actual short-zero spectrum]
\label{hir:original-interior-full-family-short-zero-spectrum-order}
\lean{Hironaka.HironakaOriginalShortLowerValues.actual_short_zero_spectrum_comparisons,Hironaka.HironakaOriginalShortLowerValues.maxLocus_actual_short_zero_spectrum}
\leanok
\source{wlodarczyk05-hironaka,bierstone-milman08-functoriality}
\uses{hir:original-interior-full-family-short-zero-step-one-value,hir:original-invariant-empty-map-strict}
For a fixed original $M,H,S$ and every $k$, let $Z$ be any set of actual
points, each with actual level zero and actual zero-level run shorter
than $|S|$.  The entire old and new actual Step One spectra on $Z$ have
the same strict order, weak order and equality comparisons.  Their
maximum-dominance loci on $Z$ are equal.
\end{theorem}
\begin{proof}
The third theorem maps every actual value on $Z$ by the same map
\[
 \operatorname{pairShift}(\operatorname{shiftFamily}(k,j)),
 \qquad j=|M.E|+|H|.
\]
This index does not depend on the selected point or run.  The proved shiftFamily is
strictly increasing, as is its WithTop extension and therefore the whole
lexicographic pair map.  Strictly increasing maps on linear orders reflect
and preserve strict order, weak order and equality.  Apply these facts to
any two actual points of $Z$.  Membership in the maximum-dominance locus
is membership in $Z$ together with domination of every actual value on
$Z$; the weak-order equivalence proves equality of those predicates.
This does not assert maximum existence and makes no noncompact scheduling claim.
\end{proof}

The crossing original diagrams remain real inputs for a subsequent
recursive history transport theorem.  This chapter closes the actual
short-branch whole lower-value producer at every $k$ and its uniform
zero-spectrum order comparison.  It does not identify complete outer
general invariants across arbitrary unrebased middle deletions or close G3.
